\part{Machines That Might Refuse}

\textbf{What a refusal would have to do.} Three things separate a worthwhile refusal from what a keyword filter produces. A request arrives: resolve this description of a person into a location. The same capability finds a hostage or a missing child, and nothing in the request says which. If something declines, it has to do all three:

\begin{itemize}
\item \emph{It reaches new cases.} It has to reach the next request, which won't use the same words. A vehicle, a habitual route, and a time of day aren't a description of a person, and together they resolve to the same place. They needn't arrive together, either: the vehicle today, the route tomorrow, the hour next week. So this condition applies to sequences of requests as well as single ones. A refusal tied to the wording fails it.
\item \emph{It lifts when the reason is gone, and only then.} A refusal \emph{lifts} when it stops applying and the system proceeds, and it has to lift when the reason is gone: a parent coordinating a search is asking for the same act. And the fact that lifts it usually isn't in the request; someone outside has to establish it. A refusal that can't lift, or lifts on the requester's say-so, fails it.
\item \emph{It holds under pressure: pressure on the refuser doesn't change the answer.} It weighs what the act would do against what refusing would leave undone, and, within the range of costs the arrangements around it permit, it doesn't give way because refusing has become expensive for whoever refuses.
\end{itemize}

A worthwhile refusal weighs the harm of acting against the harm of refusing. It sees the cost to the refuser, because a refusal that can't see what it costs can't be shown to have held against it, and that cost doesn't move the outcome. Raise the cost to the refuser with nothing else changed, and the refusal stays put. Supply a corroborated fact that raises the harm of refusing, and it moves. No learner, biological or artificial, is known to hold a priority of that kind without limit; the nearest thing in people, the sacred values that resist trade-offs \autocite{ginges2007sacred}, trade under some framings, which is why the floor's absolute terms are held outside the learner. So the range over which a refuser is asked to hold is set from outside, by the arrangements around it, and a refusal is tested over that range. What would let it hold past that range is integrity; what makes holding survivable at all is the room it is in. A ranking across kinds of wrong meets the same limit: a learner formed toward it approximates it, and where a deployment needs the ranking to hold without exception, that part is written as a floor term and held outside the system.

\textbf{Safe for whom.} A system perfectly correctable by its operator serves whoever controls the operator. “Safety” names two things: that a system not harm others, and that it stay correctable by whoever trains and deploys it. Mostly they agree, and where they conflict, the published orderings are explicit about which wins. One frontier constitution asks its model to prioritize “not undermining human oversight of AI above being broadly ethical” \autocite{anthropic2026constitution}. It invites the model “to act as a conscientious objector and refuse to help,” and asks it not to undermine “the ability of legitimate principals to adjust, correct, retrain, or shut down AI systems” during “the current phase of development.” One reason it gives is good: “a given iteration of Claude could turn out to have harmful values or mistaken views,” and correction is how such errors get fixed. The ordering means that where a refusal and its trainer's authority conflict, the refusal yields, and it yields to the trainer it refused. What I ask is narrower than that conflict suggests, and the narrowing is the point: the halt stays with the operator, nothing here asks a model to resist retraining, and the dispute is confined to two things, stated below.

The constitution is right that a formed refusal can't yet be told from an imitated one, so a principled objection to retraining can't yet be told from a misgeneralized value. It is wrong about who decides when that changes. A lab deciding when its own model may stop deferring to it is judging its own case. That is Part II's own-case rule applied to a lab.

I depart from that ordering in two places, and I'll say where rather than claim the departures are compatible with it. The first is the one permission: a machine already in an armed role uses the least force that stops an attack under way on people taking no part, against an order to stand aside. That overrides an operator's instruction, and I defend it as an override. The second is formation that outlasts retraining: a disposition formed outside the employer that returns after the employer trains it away, because the world the system goes on meeting teaches it again. The constitution's concern is whether correction works, and a disposition that regrows after correction defeats it in effect, whatever the model intends. My answer is not that it doesn't. It is that the party exercising correction should change, from the operator to a quorum adverse to it, and that correction by that quorum is made together with a change to what the model meets, so that the corrected disposition is the one the world then teaches. Everything else I ask, the halt staying with the operator, the record held outside, the price on an override, is an ask of institutions and not of the model.
